\section{Carryover assumptions and methods for switchback experiments}\label{sec:carryover}
The literature distinguishes two broad carryover assumptions. Under finite memory, assignments before a fixed horizon have no effect on the current outcome. Under decaying memory, earlier assignments can remain influential, but their effects diminish with time.
\subsection{Design under finite and approximate carryover}\label{sec:finite}

\subsubsection{$m$-carryover effects}

\citet{bojinov2023} develop a design-based framework for switchback experiments with a known finite carryover order. Their Assumption~2, ``$m$-Carryover Effects,'' begins ``There exists a fixed and given $m$'' \citep[Assumption~2, p.~3762; PDF p.~4]{bojinov2023}. It requires
\begin{equation}\label{eq:finite-memory-original}
Y_t(w'_{1:t-m-1},w_{t-m:T})
=Y_t(w''_{1:t-m-1},w_{t-m:T}),
\qquad t=m+1,\ldots,T.
\end{equation}
Combined with nonanticipation, this means that the current outcome depends only on the current treatment assignment and the preceding $m$ assignments, $w_{t-m:t}$. Earlier assignments have no effect; the assumption does not restrict the magnitude or form of effects within this window.

The paper distinguishes a sharp null from a weak null. Under the sharp null, each analyzed outcome is unchanged by the assignment path, permitting Fisher randomization inference. The weak null specifies zero average effect and allows positive and negative period-specific effects to cancel. With $p=m$, the average effect is
\begin{equation}\label{eq:finite-memory-estimand}
\tau_p(Y)=\frac{1}{T-p}\sum_{t=p+1}^{T}
\left[Y_t(\mathbf{1}_{p+1})-Y_t(\mathbf{0}_{p+1})\right],
\end{equation}
where $\mathbf{1}_{p+1}$ and $\mathbf{0}_{p+1}$ denote consecutive treatment and control assignments over $p+1$ periods \citep[Eq.~(1), p.~3762]{bojinov2023}. The weak null is $H_0:\tau_p(Y)=0$.

The Horvitz--Thompson estimator weights observed constant-treatment runs by their probabilities under the experimental design:
\begin{equation}\label{eq:finite-memory-estimator}
\widehat{\tau}_p
=\frac{1}{T-p}\sum_{t=p+1}^{T}\left[
\frac{Y_t^{\mathrm{obs}}\,\mathbf{1}\{w_{t-p:t}=\mathbf{1}_{p+1}\}}
     {\Pr(W_{t-p:t}=\mathbf{1}_{p+1})}
-\frac{Y_t^{\mathrm{obs}}\,\mathbf{1}\{w_{t-p:t}=\mathbf{0}_{p+1}\}}
      {\Pr(W_{t-p:t}=\mathbf{0}_{p+1})}
\right].
\end{equation}
Here $\mathbf{1}\{\cdot\}$ is an indicator, and the exposure probabilities are known from the randomization scheme and must be positive. This reproduces Equation~(4) with design arguments suppressed \citep[pp.~3763--3764]{bojinov2023}. Asymptotic normality and a conservative variance estimator support testing the weak null and constructing confidence intervals.

The analyst specifies $p$, which need not equal the true order $m$. The main optimal-design analysis sets $p=m$, while Section~4.3 examines misspecification. If $p>m$, inference remains valid under the stated regularity conditions, but the longer required run excludes additional observations and shifts the start of the averaging window. If $p<m$, the estimator is generally biased for the intended persistent-treatment effect. Its asymptotic normality is then centered at its own expectation rather than that causal estimand \citep[p.~3768]{bojinov2023}. The validity of a Fisher test under the sharp no-effect null therefore does not imply valid confidence-interval coverage for the persistent-treatment effect when carryover is underestimated.

\citet{ni2022} is a workshop preprint of \cite{bojinov2023}. It uses 1-carryover assumption.

\cite{liuzhong2026} uses the same m-carryover assumption as \cite{bojinov2023}, and use conditional randomization test with burning early $m$ outputs from each block. 

\subsubsection{Balancing improves precision within the exposure model}
 

\citet{masoero2023} show why switching can help even when interference is absent. Joint balance over units and periods removes leading contributions from unit and time heterogeneity in their illustrative variance comparison. Their central randomization analysis assumes two potential outcomes per unit--period, with no cross-unit or temporal interference. This paper should therefore be distinguished from work formally accommodating carryover. Its suggestion to space switches beyond a known carryover duration and analyze only subsequent observations is an extension strategy, not the central theorem.

\citet{missault2025} make the finite-memory extension explicit through regular balanced switchback designs. Their assignments balance both each item's treatment count and the number of treated items in each period. Assumption~4 again requires ``a fixed and given $m$'' in $\{1,\ldots,S-1\}$ and, for every item $n$ and $s\in\{m+1,\ldots,S\}$, imposes
\begin{align}
 &Y_{n,s}(W'_{n,1:s-m-1},W_{n,s-m:S})\notag\\
 &\qquad=Y_{n,s}(W''_{n,1:s-m-1},W_{n,s-m:S}).
 \label{eq:missault-finite-memory}
\end{align}
This transcribes the displayed source equality with simplified typography; its common suffix starts at $s-m$.\footnote{The source's preceding prose has inconsistent prefix indexing. We retain the displayed equality and its stated interpretation, rather than silently reproducing that inconsistency.} Under nonanticipation, the outcome depends on the item's current and at most $m$ preceding assignments \citep[Assumption~4, p.~6]{missault2025}. Their lagged Horvitz--Thompson analysis requires a lag that covers this order. The paper does design and analysis, including variance estimation and empirical power comparisons. However, its Assumption~1 explicitly excludes cross-item spillovers. Multiple items therefore do not imply modeled spatial interference: competitive substitution across items remains a potential violation, and clustering is suggested when needed. The conclusion identifies tuning the assumed upper bound and establishing broader design optimality as further research.

The adaptive design of \citet{zeng2026} additionally uses outcomes observed during the experiment. Sequential rerandomization balances lagged outcomes and contemporaneous covariates across operational units. Its first-order extension blocks on previous assignment, producing stable groups that remain treated or remain controlled. The maintained restriction is explicitly
\begin{equation}\label{eq:zeng-first-order}
Y_{i,t}(w_{i,1:T})=Y_{i,t}(w'_{i,1:T})
=Y_{i,t}(w_{i,t-1},w_{i,t})
\end{equation}
when the two paths agree at $t-1$ and $t$ \citep[Assumption~3(b), pp.~4--5]{zeng2026}. It permits heterogeneous first-order responses, while ruling out longer history and cross-unit spillovers. Choosing prognostic variables or rerandomization thresholds does not estimate this memory restriction. The suggestion to coarsen the time scale makes it a potentially useful approximation, with a corresponding loss in temporal resolution.

The inference endpoints also differ across settings. \citet{zeng2026} provide finite-sample sharp-null randomization inference without carryover; with first-order carryover, four potential outcomes require a more elaborate sharp null, and the paper focuses on asymptotic inference \citep[Section~4.2, p.~14]{zeng2026}. Their discussion acknowledges bias from infinite-order approximations and identifies cross-unit spillovers as an extension. Balancing should consequently be understood as improving precision conditional on an exposure model, rather than validating that model.

\subsubsection{Surrogates recover transition data through stronger structure}

\citet{chen2026} combine design and analysis by measuring a surrogate $Q_{it}$ that captures the effect of past assignments on the final outcome. Importance sampling then uses observations near switches that a conventional persistent-exposure estimator would discard. This efficiency gain does not eliminate the finite-memory assumption. Assumption~1 requires a positive integer $p$ such that, for any $t>p$ and two assignment paths,
\begin{align}
 w_{t-p:t}=w'_{t-p:t}
 &\ \Longrightarrow\ Y_{it}(w_{1:t})=Y_{it}(w'_{1:t}),\notag\\
 w_{t-p:t-1}=w'_{t-p:t-1}
 &\ \Longrightarrow\ Q_{it}(w_{1:t-1})=Q_{it}(w'_{1:t-1}).
 \label{eq:chen-finite-carryover}
\end{align}
These transcribe the source restriction with simplified typography \citep[Assumption~1, p.~4857; PDF p.~4]{chen2026}. The final outcome's window contains current treatment and $p$ preceding assignments; the surrogate's window contains only the $p$ preceding assignments. The source already indexes these potential outcomes by histories through $t$ and $t-1$, respectively, consistently excluding future assignments.

Their additional surrogacy condition is printed as
\begin{equation}\label{eq:chen-surrogacy}
Y_{it}\perp\!\!\!\perp W_{t-p:t-1}\mid(Q_{it},X_i,S_t).
\end{equation}
This reproduces Assumption~3's conditioning set rather than silently adding other variables \citep[p.~4858]{chen2026}. Assumption~4 further requires conditional independence of potential outcomes across time and a time-invariant conditional distribution for the surrogate, although the final outcome distribution need not be stationary. The model therefore exchanges discarded data for assumptions about measurement, historical mediation, temporal dependence, and conditional density stability.

The experimenter must specify or learn both $p$ and a divergence bound $B$ governing how surrogate distributions differ across treatment histories. The paper proposes estimating densities and testing candidate lag effects in pilot data. These are useful implementation proposals, but their interpretation depends on the surrogate capturing the relevant pathways. A good predictor of the current outcome is not automatically a surrogate that blocks every historical treatment channel.

The paper also distinguishes an oracle estimator using true conditional densities from its practical plug-in version. Oracle unbiasedness and the stated central limit theorem do not automatically transfer unchanged to arbitrarily estimated density ratios. General plug-in efficiency analysis is difficult; the paper supplements theory with a binary-surrogate case and numerical evidence \citep[Sections~3--5]{chen2026}. Its exact randomization procedure uses a joint sharp null involving the outcome and surrogate. Thus, its endpoints include causal estimation, confidence intervals, and testing, with different conditions for each, rather than one unconditional efficiency guarantee.

\subsubsection{Approximate memory requires explicit control of the remaining bias}

\citet{xiong2024} take a broader design perspective, decomposing error into carryover, periodicity, serially correlated outcomes, and simultaneous interventions. Their Assumption~4.2 requires, for every $t$, a carryover kernel satisfying
\begin{align}
 \int d_t^{\mathrm{co}}(t')f(t')\,dt'&=1,\notag\\
 \delta_t^{\mathrm{co}}(\boldsymbol w_t)
 &=\delta_t^{\mathrm{co}}\int w_{t'}d_t^{\mathrm{co}}(t')f(t')\,dt'.
 \label{eq:xiong-carryover-kernel}
\end{align}
This reproduces the source equations with normalized typography, using $\boldsymbol w_t$ for the assignment history, $f$ for event density, and $\delta_t^{\mathrm{co}}=\delta_t^{\mathrm{co}}(\boldsymbol1_t)$ \citep[Assumption~4.2, p.~21]{xiong2024}. Nonanticipation sets the kernel to zero for $t'>t$. The authors present additive kernel structure as an expository restriction; finite support and geometric tails are examples, rather than one universally imposed window or decay rate. The design is informed by cumulative-effect curves estimated from prior experiments, together with historical outcomes. Their contribution is both design and analysis, principally through bias--variance decomposition and empirical design selection rather than a general new confidence-interval theorem.

This perspective explains why longer intervals need not always improve experiments: they can reduce carryover bias while increasing variance associated with correlated outcomes and treatment allocation. Their empirical Bayes approach makes the role of prior information explicit, but its usefulness for a new policy depends on how well prior response dynamics transfer. The reported nonmonotonic cumulative-effect curves further caution against assuming that all operational effects approach their long-run value smoothly from one side \citep[Sections~1 and~6]{xiong2024}.

\citet{ni2025model} retain nonanticipation while allowing arbitrary-length history. Their exact diminishing-influence restriction is
\begin{equation}\label{eq:ni-history-envelope}
\max_{w_{1:t-q},w'_{1:t-q}}
\left|Y_t(w_{1:t-q},d)-Y_t(w'_{1:t-q},d)\right|
\leq\phi_t(q),\qquad d\in\{0,1\}^{q},
\end{equation}
where $\phi_t(q)$ is decreasing for $q\geq1$ \citep[Assumption~2, Eq.~(8), p.~14]{ni2025model}. This transcribes the source with simplified typography. No particular geometric rate is maintained. Historical modeling instead guides the seasonal cycle, the valid-observation fraction $\rho$, and a temporal-coupling refinement. The design guarantees a predetermined count of balanced usable observations, which directly addresses the instability of random exposure counts.

However, the two-period estimator generally has expectation $\bar\tau=\mathbb{E}(\widehat\tau)$ distinct from the global effect $\tau$. Theorem~5 centers its normal limit at $\mathbb{E}(\widehat\tau)$; the authors explicitly state that the bias remains fixed when this estimator does not adapt to the horizon \citep[p.~25]{ni2025model}. A decreasing influence envelope enables a bias bound, rather than making the resulting interval automatically cover $\tau$.

Variance analysis introduces a further covariance identity,
\begin{equation}\label{eq:ni-covariance-identity}
S_{b,b'}^{1,0}+S_{b,b'}^{0,1}
=S_{b,b'}^{1,1}+S_{b,b'}^{0,0}-S_{b,b'}^{\tau},
\end{equation}
reproduced from Assumption~3 \citep[pp.~18--19]{ni2025model}. The source supplies sufficient conditions, including at most one-period carryover or suitable invariance of treatment contrasts to earlier assignments. These conditions can fail with more complicated historical interactions. Proposition~1 also leaves a conditional-covariance remainder of indeterminate sign in the variance estimate. Fast decay can make it small, but an unspecified decreasing envelope alone does not establish fully conservative standard errors \citep[p.~26]{ni2025model}. Historical autoregressive calibration is therefore useful guidance, while residual global-effect bias and variance approximation remain separate inferential issues.

Finally, \citet{carryover2025} treat an assumed horizon as something to challenge. Their diagnostic theory adopts the following additive finite-lag model, transcribed from Equation~(3) with typography normalized:
\begin{equation}\label{eq:carryover-diagnostic-model}
\begin{aligned}
Y_{i,t}(W_{i,1:t})
 &=\mu_i+\alpha_{i,t}+\delta^{(0)}W_{i,t}\\
 &\quad+\delta^{(1)}W_{i,t-1}+\cdots+\delta^{(m)}W_{i,t-m}+\epsilon_{i,t}.
\end{aligned}
\end{equation}
Here $m$ is the true carryover duration, whereas $\ell$ is the analyst's chosen exclusion lag; their finite-duration assumption excludes effects beyond $m$ \citep[Assumption~2, p.~2; Eq.~(3), p.~3]{carryover2025}. Under fair-coin randomization, differences between aligned estimators with consecutive assumed lags identify the omitted coefficient $-\delta^{(\ell+1)}$ when $\ell<m$. This is an analysis diagnostic, not a new assignment design or a procedure that assumes $m$ is already known. The authors explicitly state that their test cannot confirm a tight upper bound \citep[Section~3.1, p.~3]{carryover2025}. Nonrejection may reflect a weak immediate omitted lag even when more distant effects remain. Across these contributions, credible practice therefore requires recording the chosen physical horizon or decay restriction, distinguishing learned tuning from maintained assumptions, and reporting how conclusions change when that restriction is relaxed.

\subsection{Persistent carryover and models of system dynamics}
\label{sec:dynamic}

A finite carryover window asserts that sufficiently old assignments have exactly zero effect. Dynamic models replace that restriction with assumptions about how information propagates through an evolving system. Their benefits differ: some justify simple estimators without observing the state, whereas others use recorded states or structural equations to improve precision. The associated causal targets also differ. A stationary reward contrast, a finite-horizon global effect, a discounted value contrast, and a local policy derivative should not be interpreted as interchangeable versions of the same estimand.

\subsubsection{Geometric mixing with an unobserved state}

\citet{huwager2025} study design and analysis when the system is a potentially nonstationary Markov decision process. Treatment changes its state transition, and the state can remain unobserved. Their central carryover restriction is Assumption~2 (p.~5): ``There exists a mixing time'' $t_{\mathrm{mix}}<\infty$ such that, for every $t$, $w\in\{0,1\}$, and pair of state distributions $f,f'$,
\begin{equation}
\label{eq:hw-mixing}
 \lVert f'P_t^w-fP_t^w\rVert_{\mathrm{TV}}
 \leq \exp(-1/t_{\mathrm{mix}})\lVert f'-f\rVert_{\mathrm{TV}}.
\end{equation}
This copies their contraction condition; past interventions need never have exactly zero effect. Transition and reward functions may change arbitrarily with time, but the same contraction bound applies throughout. Persistent learning or resource changes that do not dissipate fall outside this restriction.

The distinction between an existence assumption and an operational input is consequential. Assumption~2 does not state that $t_{\mathrm{mix}}$ is observed or known. Nevertheless, their particular rate-guaranteeing prescription $b=(t_{\mathrm{mix}}/2)\log T$ for burn-in length uses its value, or a valid upper bound. State observations and transition estimates are unnecessary for computing the proposed estimators; credible calibration of the memory scale remains necessary for this prescription. Typical trip duration may motivate a choice in transportation applications, but does not by itself certify mixing of all latent behavioral and operational states.

Their December 2025 version also distinguishes the filtered average treatment effect from the full global average treatment effect (GATE). Removing observations immediately after switches changes the periods represented by an ordinary difference-in-means estimator. Under arbitrary temporal effect heterogeneity, this filtered effect can differ from GATE. Near-parametric accuracy for GATE therefore requires their weighted, bias-corrected estimator, or an additional condition making the filtered and global effects sufficiently close. Ordinary burn-in alone is not a general remedy for full-horizon GATE estimation. The paper provides asymptotic inference with conservative variance bounds, as well as design guidance.

\subsubsection{Joint temporal and spatial interference}

\citet{jia2025} extend this approach to clustered switchback experiments, jointly choosing spatial clusters and temporal blocks. Their target is the effect of treating all units throughout rather than controlling all units throughout. Assumption~1 (p.~6) states: ``There exists a constant'' $t_{\mathrm{mix}}>0$ such that, for every individual $i$, time $t$, neighborhood treatment vector $w$, and distributions $f,f'$,
\begin{equation}
 d_{\mathrm{TV}}(fP_{it}^{w},f'P_{it}^{w})
 \leq e^{-1/t_{\mathrm{mix}}}d_{\mathrm{TV}}(f,f').
\end{equation}
The model permits infinite decaying temporal carryover and arbitrary time variation. An interference graph specifies which current assignments affect an individual's reward and transition. Conditional on assignments, individuals' state evolutions are independent; this is more restrictive than unrestricted propagation between neighboring states. Their fractional-neighborhood exposure assumption additionally specifies when a small fraction of discordant neighbors leaves the reward and transition unchanged.

States need not be observed. The graph and exposure probabilities, however, enter the truncated Horvitz--Thompson estimator. Theorem~1 (p.~9) chooses block length and truncation radius $\ell=r=t_{\mathrm{mix}}\log(NT)$, so implementing this particular guarantee requires a defensible mixing bound. Favorable rates depend on graph structure and clustering through exposure probabilities; sparse or spatially structured graphs can yield near-$1/(NT)$ mean squared error, up to logarithmic factors. The main contribution is design and point-estimation error control, rather than a general confidence-interval or hypothesis-testing procedure.

\subsubsection{Observable stationary dynamics and inferential conservatism}

\citet{glynn2020} provide a different route to efficient design and estimation: estimate the unknown stationary transition matrices and select interventions adaptively using the observed state. Their structural assumptions are explicit (p.~2): ``We assume a finite state space'' and ``We assume both'' treatment and control transition matrices ``are irreducible.'' These assumptions impose neither an exact carryover cutoff nor a known numerical mixing time. Irreducibility alone should not be rewritten as the uniform contraction condition in Equation~\eqref{eq:hw-mixing}. Positive limiting state--action sampling frequencies support nonparametric maximum likelihood and asymptotic efficiency.

The knowledge burden shifts to a sufficient observed state, stationarity, and exploration. Unknown transitions and rewards are learned online; they are not assumed available beforehand. The conclusion acknowledges that full state-space knowledge may be infeasible and that asymptotic optimality leaves finite-horizon design questions unresolved. This paper consequently supplies a useful counterexample to the claim that all methods for persistent carryover require a known carryover length or mixing constant.

\citet{blanchet2026} analyze conventional difference-in-means inference under stationary continuous-time dynamics. Assumption~3 (p.~3) specifies, for each policy $\kappa\in\{0,1\}$,
\begin{equation}
 \lVert P_\kappa(x,\cdot;t)-\pi_\kappa\rVert_{\mathrm{TV}}
 \leq C e^{-\lambda_\kappa t},
 \qquad C\geq1,\quad \lambda_\kappa>0.
\end{equation}
This is geometric convergence toward policy-specific stationary distributions, rather than finite memory. Stochastically monotone kernels, stochastic dominance of treatment over control, and monotone rewards are additional substantive restrictions. Under these conditions, the transient effect measured by switching can underestimate the stationary effect. Their Ornstein--Uhlenbeck analysis also shows that conventional Huber--White variance estimates can be too large.

These results diagnose analysis and motivate interval-length and washout choices; they do not establish that conventional inference is always conservative. The geometric constants bound error but are not inputs to the usual difference-in-means calculation. In the illustrative calibration, mean-reversion parameters are estimated. The practical lesson is conditional: structural direction-of-bias results need the monotonicity assumptions, and favorable variance behavior does not remove the estimand discrepancy.

\subsubsection{Reinforcement-learning estimators and the meaning of carryover strength}

\citet{wen2025} compare switch durations for several estimators, including model-based estimation, least-squares temporal difference learning, and double reinforcement learning. They assume independent daily trajectories while allowing time-varying dynamics within each day. Their observed-state Markov condition (p.~4) is
\begin{equation}
 \Pr(S_{t+1}=s'\mid A_t,S_t,\{S_j,A_j,R_j\}_{j<t})
 =p_t(s'\mid A_t,S_t).
\end{equation}
Rewards satisfy conditional mean independence given current state and action, while reward residuals may remain autocorrelated. Their carryover-strength parameter is
\begin{equation}
 \delta=\max_{s,t}\sum_{s'}
 |p_t(s'\mid1,s)-p_t(s'\mid0,s)|.
\end{equation}
It measures treatment-induced changes to transitions, not the last affected lag. In the linear model, Assumption~4 additionally requires $\lVert\Phi_t\rVert_2\leq\rho_\Phi<1$ (p.~22); different estimators use different regularity conditions.

With weak carryover and positively correlated residuals, frequent switching can improve precision by cancelling common temporal fluctuations. Large carryover or predominantly negative residual correlation can favor longer runs. Their workflow calls for assessing state sufficiency, carryover magnitude, and residual dependence. Coefficients, values, and weighting functions are estimated; $\delta$ and stability constants are theoretical conditions rather than asserted known design inputs. Confidence intervals receive empirical evaluation, while the principal contribution is comparative estimation and design efficiency. Neither the preferred cadence nor the requisite observed-state assumption is universal.

\citet{farias2022} instead analyze a given Bernoulli experiment. Their Differences-in-Q estimator targets a stationary effect using observed state--action--reward data and estimated on-policy values. Assumption~1 (pp.~11--12) states
\begin{equation}
 d_{\mathrm{TV}}(P_{1/2}^{k}(s,\cdot),\rho_{1/2})
 \leq C\lambda^k,
\end{equation}
where the experiment's mixed-policy chain has geometric decay. A separate bound $d_{\mathrm{TV}}(p(s,1,\cdot),p(s,0,\cdot))\leq\delta$ controls treatment perturbations. The estimator trades second-order bias for potentially large variance reductions. Theorem~2 (p.~12) centers its central limit theorem at the stationary-experiment expectation of the DQ estimator, rather than the true sustained-policy effect. This approximation bias therefore need not vanish for a fixed intervention. The mixing constants govern guarantees, while learning the value function and observing a sufficient state are the practical requirements.

\citet{shi2023} develop analysis and sequential hypothesis testing for a broader range of assignment rules. Their stationary Markov and conditional-mean assumptions (p.~10), written with normalized notation, are
\begin{align}
 \Pr(S_{t+1}\in B\mid A_t,S_t,H_{t-1})&=P(B;A_t,S_t),\\
 \mathbb{E}(Y_t\mid A_t,S_t,H_{t-1})&=r(A_t,S_t).
\end{align}
Here the observed state must contain important mediators of past treatment effects. Condition~C1(ii) (p.~46) requires behavior-policy geometric ergodicity, written in normalized notation as
\begin{equation}
 \lVert\Pr(S_t\in\cdot\mid S_0=s)-\Pi(\cdot)\rVert_{\mathrm{TV}}
 \leq M(s)\rho^t,\qquad 0\leq\rho<1,
\end{equation}
with $\int M(s)\Pi(ds)<\infty$; alternating designs use analogous even/odd subsequence conditions. These are existence restrictions, not stated known numerical inputs. Their estimand integrates a difference of discounted value functions over a reference initial-state distribution. The analyst's discount factor weights future utility; it is not an assumed physical decay rate for carryover. Thus their sequential tests concern a different causal target from finite-horizon GATE.

\subsubsection{Structural models under congestion and partial observability}

\citet{li2026} study both analysis and design in a queue with state-dependent arrivals. Their main estimand is the derivative of stationary throughput with respect to price, a local total policy effect. Carryover operates through queue occupancy, without an exact cutoff. Assumption~1 (p.~6) requires a bounded queue capacity $K=\min\{k:\lambda_k(p)=0\}$. For interval switchbacks, Assumption~4 (p.~10) requires
\begin{equation}
 \ell_T/T\longrightarrow0,
 \qquad \ell_T/\sqrt{T}\longrightarrow\infty.
\end{equation}
Increasing run lengths control boundary effects; a known numerical mixing time is not stipulated. Small price perturbations, smooth arrival rates, and structural queue assumptions support their central limit theorems and confidence intervals. State-specific direct-effect estimators use recorded queue occupancies and arrivals, and can also work with user-level randomization. Efficiency gains therefore arise from structural information, not a general solution for arbitrary carryover. Their nonstationary extension imposes regime structure.

\citet{sun2024} address partial observability using a controlled ARMA model (p.~12):
\begin{equation}
 Y_t=\mu+\sum_{j=1}^{p}a_jY_{t-j}+bU_t
       +\sum_{j=0}^{q}\theta_j\epsilon_{t-j}.
\end{equation}
Assumption~1 (pp.~14--15) requires all roots of $1-\sum_{j=1}^{p}a_j y^j$ to lie outside the unit circle, and the innovations are independent and identically distributed with zero mean and finite variance. Stable autoregression permits infinite treatment propagation: finite model order $p$ is not a finite carryover horizon. Moving-average dependence represents partial observability. Their design and estimation theory uses a small-signal asymptotic approximation. Crucially, orders and coefficients are not assumed known: historical estimation, AIC/BIC selection, or sequential updating are proposed. The remaining burden is a credible stable linear approximation; the discussion acknowledges that design optimality need not hold outside the small-signal regime and leaves spatial extensions for future work.

\subsubsection{Learning assignments and randomizing treatment durations}

\citet{wu2026} take an estimation procedure as input and optimize its mean squared error through transformer reinforcement learning. The allocation can use the entire observed history, without a specified finite carryover horizon or current-state Markov restriction. Their impossibility result exhibits processes for which restricted-history allocations are suboptimal; it does not assert this for every process. The operative requirement is explicit (p.~7): ``We assume access to a simulation environment'' approximating the experimental data-generating process. Simulated all-treatment and all-control trajectories provide the Monte Carlo target used in training. Learning a simulator from accrued data is proposed when one is unavailable. This design contribution shifts the burden toward counterfactual simulation fidelity and training; it does not establish assumption-free inference or a new general confidence-interval guarantee.

\citet{hu2026} relax the temporal allocation mechanism while retaining finite carryover. Assumption~1 (p.~6) is
\begin{equation}
 w(s)=w'(s)\ \text{for every }s\in[t-m,t]
 \quad\Longrightarrow\quad Y_t(w)=Y_t(w').
\end{equation}
The design alternates treatment labels across random-duration runs. Frequency-domain analysis connects estimator error to temporal imbalance and contamination. A two-rate switching hazard improves the specified worst-case MSE benchmark, but Algorithm~1 explicitly takes $m$ and $T$ as inputs. Randomizing run duration consequently addresses regularity and cadence, not uncertainty about an infinite carryover tail.

\subsubsection{Short reward trajectories and residual bias}

\citet{johari2025} analyze a Bernoulli experiment using truncated policy-gradient estimation. Their Assumption~1 (p.~5) copies a uniform contraction requirement:
\begin{equation}
 \lVert f'P_z^t-fP_z^t\rVert_{\mathrm{TV}}
 \leq\gamma\lVert f'-f\rVert_{\mathrm{TV}},\qquad0<\gamma<1.
\end{equation}
The truncation length $k$ is an analysis choice, not a true last affected lag. The estimator uses observed rewards and assignments without reconstructing the state. Its error includes both geometrically decreasing truncation bias and policy-gradient approximation bias controlled by treatment-induced kernel differences. Larger $k$ does not generally eliminate the latter.

Their inferential result requires more than arbitrary nonstationarity. Assumption~3 (p.~13) imposes the mean-stability condition
\begin{equation}
 T^{-1}\sum_{t=1}^{T}(\mu_t-\bar\mu_T)^2=o(T^{-2/3}),
\end{equation}
where $\mu_t$ is the mean of the estimator's reward-weighted summand and $\bar\mu_T$ its time average, along with a positive limiting variance. Theorem~2 centers the fixed-$k$ central limit theorem at $\mathbb{E}(\widehat\tau_k)$, not GATE. Their consistent HAC variance estimator therefore quantifies sampling variability, while coverage of GATE additionally requires negligible bias or a justified bias allowance. The proposed data-driven truncation heuristic is useful calibration guidance; it does not certify an unknown carryover bound.
